\section{百篇精读后优先推进的六个问题}
下面的“未解决”均指所读证据尚未给出跨模型、跨参数的充分答案，不声称全世界没有其他研究。每项同时说明Hall/双流体能够推进到哪一步。

\begin{longtable}{P{30mm}P{57mm}P{61mm}}
\toprule 问题 & 已有证据与剩余缺口 & 可执行的下一步\\\midrule\endhead
电子层的真实厚度与慢排出 & MRX约$8d_e$、$0.11V_{\!Ae}$；经验摩擦改善但未闭合压力/涨落预算\citep{L119} & 同装置边界、同探针卷积，比较五矩、十矩、非局地热流；同时匹配宽度、速度、$E$残差，不能只调一条总率曲线\\
三维斜模如何成为非线性X线 & 惯性force-free与电阻Harris线性结果不同；倾角误差不足以区分两选择准则\citep{L224,L268,L166} & 固定同一平衡，扫描电子非理想项与允许波数，比较完整三维模式竞争；区分相速度和前沿速度\\
爆发加速何时被岛/涡旋截断 & 冷模型增加横向分辨率后出现岛；EMHD岛内射流另有KH类失稳\citep{L191,F697,E459} & 幅度定义统一，网格与滤波双扫描；诊断局域X几何、岛出现时间、能量守恒，勿只拟合一个幂指数\\
传播记忆能否在真实磁尾保持 & 非均匀片厚使前沿依赖起始片厚；当前冷电子小质量比模型有限\citep{L294} & 在同一几何加入电子压强/导场和不同载流份额，检验历史效应是否仍存在；装置测量加入有限时间/空间分辨\\
高阶矩能否自洽预测粒子分布 & 最大熵重建正定但21矩仍缺多峰；重建尚非时演验证\citep{E553} & 先检验可实现性和守恒，再从相同初值演化；同时比较压力、热流、分布距离与重联率，而非只拟合已知PIC矩\\
磁能转移如何变成不可逆热 & Hamiltonian能量转移无正定熵源；KREHM给条件性相混合热化机制\citep{E512} & 独立扫描空间与Hermite分辨率、碰撞频率和观察时间；区分数值耗散、压力功和碰撞熵增\\
\bottomrule
\end{longtable}

\subsection{哪些边界不能靠更细Hall网格消除}
标量压强Hall模型没有任意电子非旋转对称压力、俘获/反向粒子群、速度空间相混合和粒子高能尾的自由度。
增加网格能减少离散误差，却不会使方程长出这些未知量。
十矩能提供压力张量通道，但热流仍需闭合；最大熵矩族能增加分布形状的表达能力，却有可实现性和计算代价问题。
KREHM保留电子平行动理学，但其低$\beta$强导场排序又排除了一般反平行高$\beta$情形。
因此一个模型不能仅因在GEM总磁通上成功，就被认为同时解决了磁尾分布、实验加热和高导场湍流。

\subsection{几项值得研究但不能过早宣称的推断}
\citet{F697}中的太阳耀斑时间应用依赖很强的参数外推，部分数值接近相对论或不满足强导场序；需先确认排序，再讨论观测意义。
\citet{L313}关于指数拉伸和电子惯性的框架也需给定具体三维流及边界，才有可验证的时间预测。
\citet{F644}的无显式电子耗散Hall湍流应先量出离散电场残差，才能判断小片重联是否在收敛极限存在。
这些不是否定研究方向，而是给后续模拟设置明确的成功标准。
